% ARCHIVED AUDITED DRAFT: independently checked by root, 2026-11-19 15:00 UTC.
% Incorporated into part_B.tex, whose version is authoritative.
\begin{saresult}{B.051}{Proved}{Completion in every Noetherian dimension}
Let $(R,\mathfrak m)$ be a Noetherian local domain with
fraction field $K$, let $A=\widehat R$, $Q=K\otimes_R A$,
and let $R\subset L\subset K$ be any overring.
Put $B=L\otimes_R A\subset Q$. For every $n$,
\[
                    \Omega_n(R,L)\otimes_R A
                         \simeq\Omega_n(A,B).              \tag{B51.1}
\]
The tensor maps and existence of the two polynomial
co-operations pass through this identification. There is
no dimension or finite-normalization hypothesis.
\end{saresult}
\begin{proof}
If the residue field is infinite, a grid of points with
distinct residues gives an invertible Vandermonde
matrix over both local rings, so the integer-valued
modules are the ordinary polynomial rings. Assume
the residue field is finite.

For $j\ge0$ define
\[
 C_j=\{x\in K:\mathfrak m^jx\subset R\},\qquad
 C'_j=\{x\in Q:(\mathfrak m A)^jx\subset A\}.
\]
Flatness gives $C_j\otimes_R A=C'_j$.
Indeed, choose finitely many generators $a_1,\ldots,a_s$
of $\mathfrak m^j$ and tensor the kernel description
\[
 C_j=\ker\bigl(K\longrightarrow(K/R)^s,
                       x\longmapsto(a_ix\bmod R)_i\bigr).
\]

Every coefficient of an integer-valued polynomial in
$K[X_1,\ldots,X_n]$ lies in some $C_j$, by the
punctured-spectrum Vandermonde argument of B.034.
We verify the corresponding assertion over $A$ even
though $A$ need not be a domain and $Q$ need not be
Artinian. Each coefficient $x\in Q$ has a denominator
$d\in R\setminus\{0\}$, which is a nonzerodivisor
of $A$ by flatness. If $\mathfrak p$ is a nonmaximal
prime of $A$, its contraction to $R$ is not
$\mathfrak m$, since otherwise it contains
$\mathfrak m A$, the maximal ideal of $A$.
The image of $R$ modulo this contraction is infinite.
Choose a rectangular Vandermonde grid of $R$-points
whose determinant $v$ is outside $\mathfrak p$.
The values of the polynomial are in $A$, so Cramer's
rule gives $vx\in A$. Consequently the ideal
\[
                         I_x=\{a\in A:ax\in A\}
\]
is not contained in any nonmaximal prime. If it is
proper, its radical is the maximal ideal, and finite
generation gives $(\mathfrak m A)^j\subset I_x$
for some $j$. Thus $x\in C'_j$.
Taking a common $j$ handles all coefficients of a
given polynomial.

Fix this $j$ and a finite rectangular degree bound
with $N$ allowed monomials. Let $E_j\subset R$
represent the finite ring $R/\mathfrak m^j$.
For coefficients in $C_j$, the value modulo $R$
depends only on the inputs modulo $\mathfrak m^j$:
a monomial difference has a factor in that ideal,
which sends each coefficient into $R$.
The desired bounded module is therefore
\[
 \ker\bigl(C_j^N\longrightarrow(C_j/R)^{E_j^n}\bigr).
\]
Since $R/\mathfrak m^j\simeq A/(\mathfrak m A)^j$,
the same finite tests apply over $A$. Tensoring the
kernel gives exactly the bounded module there.
Take the filtered union over $j$ and the degree
bounds to prove (B51.1) for $L=K$.
For arbitrary $L$, intersect the bounded coefficient
module with $L^N$; flatness preserves this finite
intersection, as in B.047. Tensor compatibility and
existence follow from multiplication and finite
$A$-linear combination exactly as in B.047.
\end{proof}

\begin{saresult}{B.052}{Proved}{Epimorphic enlargements supported at the closed point}
Let $(A,\mathfrak m)$ be a complete Noetherian local
ring, containing a complete discrete valuation ring
$V$ with uniformizer $t$, with $t\in\mathfrak m$
a nonzerodivisor on $A$. Assume the residue field
of $V$ maps isomorphically to $A/\mathfrak m$.
Let $W\subset A$ be a multiplicative set of
nonzerodivisors containing $t$, and put $Q=W^{-1}A$.
Suppose
\[
 A\subset B\subset Q,\qquad
 \text{every element of }B/A\text{ is killed by a power
 of }\mathfrak m.
\]
Then there is a finite-dimensional
$F=\operatorname{Frac}(V)$ vector space $D\subset B$
which is an ideal of $Q$, such that $B/D$ is a finite
$A$-module and $A\cap D$ is a finite $V$-module.
If $A\to B$ is a ring epimorphism, then
\[
                         B=A+eQ,\qquad D=eQ                 \tag{B52.1}
\]
for an idempotent $e\in Q$, and $eQ$ is finite-dimensional
over $F$. If the common residue field is finite,
injectivity of the polynomial tensor map for
$\Omega(A,B)$ implies this epimorphism condition.
\end{saresult}
\begin{proof}
Put $T=B/A$. Since every class is killed by a power
of $\mathfrak m$, it is killed by a power of $t$,
and $B\subset A[1/t]$.
The $t$-torsion subgroup $T[t]$ embeds in $A/tA$:
send the class of $a/t$ to $a\bmod tA$.
Its image is contained in the submodule of $A/tA$
killed elementwise by powers of $\mathfrak m$.
The ascending chain of annihilator submodules
$(0:_{A/tA}\mathfrak m^j)$ stabilizes, since
$A/tA$ is Noetherian. The stable module is killed
by one power of $\mathfrak m$, and hence has finite
length. In particular $T[t]$ is a finite-dimensional
vector space over $V/tV$.

Let $E=F/V$. By the injectivity of $E$ proved in
B.048, an embedding $T[t]\hookrightarrow E[t]^s$
extends to a map $T\to E^s$. This extension is
injective: any nonzero kernel element, being killed
by a power of $t$, has a nonzero multiple killed
by $t$, contradicting injectivity on $T[t]$.
The same double-annihilator argument as in B.048
classifies submodules of $E^s$. Explicitly, annihilate
$T$ by a submodule of
$\operatorname{Hom}_V(E^s,E)=V^s$; cogeneration by
$E$ shows that taking the annihilator again gives
$T$. Smith reduction of that finite submodule of
$V^s$ shows
\[
                            T\simeq E^r\oplus T_0,
\]
where $T_0$ has finite length over $V$.
Its maximal divisible submodule $T_{\mathrm{div}}$
is $E^r$ and is canonical, hence stable under $A$.

The ring $A$ is also $t$-adically complete and
separated. To see completeness, a $t$-adic Cauchy
sequence is $\mathfrak m$-adically Cauchy; its limit
has the required congruences because each ideal
$t^jA$ is $\mathfrak m$-adically closed.
Closedness and separatedness are the Krull
intersection theorem applied to finite modules.
There is consequently a natural isomorphism
\[
 \operatorname{Hom}_V(E,A[1/t]/A)\simeq A.                  \tag{B52.2}
\]
Indeed the image of $t^{-j}+V$ is
$a_j/t^j+A$ with $a_j\in A/t^jA$; the homomorphism
condition is the compatibility of these residues,
and completeness gives a unique $a\in A$.

Choose the $r$ coordinate injections of
$T_{\mathrm{div}}\simeq E^r$ into $A[1/t]/A$.
By (B52.2) they are multiplication by elements
$a_1,\ldots,a_r\in A$. Put
$D=Fa_1+\cdots+Fa_r$.
These vectors are $F$-linearly independent:
a relation, after scaling to coefficients in
$V$ with at least one unit, would give a nonzero
kernel on $E^r$.
Their fractional multiples lie in $B$ because
their classes lie in $T$. Moreover
\[
              A\cap D=Va_1+\cdots+Va_r,\qquad
              D/(A\cap D)=T_{\mathrm{div}}.
\]
The first equality follows from injectivity of
the coordinate map $E^r\to T$.
Since $B/(A+D)\simeq T_0$ is finite over $V$,
$B/D$ is a finite $A$-module.

Krull intersection on this finite module shows
$\bigcap_j t^j B=D$. Thus $D$ is an ideal of
$B$, and in particular stable under $A$.
Every $w\in W$ acts injectively on the finite
$F$-vector space $D$, because it is a unit of $Q$.
It therefore acts bijectively. Hence $D$ is stable
under $W^{-1}A=Q$, so it is an ideal of $Q$.

Suppose now $A\to B$ is a ring epimorphism.
The finite $A$-algebra $B/D$ is an epimorphic image
of $A$. If it is zero the map from $A$ is surjective;
otherwise the residue-field dimension argument
and Nakayama proof of B.049 show that $A\to B/D$
is surjective. Thus $B=A+D$. With $I=A\cap D$,
the same tensor computation gives
\[
 B\otimes_A(B/A)=0,\qquad B\otimes_A D=D,\qquad
                         D=BI=I+DI.
\]
Here $B\otimes_A D=D$ follows by localizing at
$W$ on the already $Q$-module $D$.
The quotient $D/DI$ is an $F$-vector space generated
over $V$ by the finite module $I$, so it vanishes.
Therefore $D=D^2$.
Since $D$ is finite-dimensional over $F\subset Q$,
it is finitely generated as an ideal of $Q$.
A finitely generated idempotent ideal is generated
by an idempotent: write generators $d_i$ as
$d_i=\sum a_{ij}d_j$ with $a_{ij}\in D$;
the adjugate identity says $\det(1-(a_{ij}))D=0$.
Write this determinant as $1-e$, with $e\in D$.
Then $(1-e)e=0$, so $e^2=e$, and $D=eQ$.
This proves (B52.1).

Finally suppose the residue field is finite and
the polynomial tensor map is injective. For any
$b\in B$ choose $j$ with $\mathfrak m^j b\subset A$.
A monic null polynomial of the finite ring
$A/\mathfrak m^j$ gives $P,bP\in\Omega(A,B)$.
Top-coefficient extraction from
$bP\otimes P-P\otimes bP$ gives
$b\otimes1=1\otimes b$ exactly as in B.049.
Thus $B\otimes_A B\simeq B$.
\end{proof}

\begin{saresult}{B.053}{Proved}{The higher-dimensional local classification without finite normalization}
Let $(R,\mathfrak m)$ be a Noetherian local domain
of dimension at least two with finite residue field,
let $K=\operatorname{Frac}(R)$ and $R\subset L\subset K$.
Put
\[
 R^\sharp=\{x\in K:\mathfrak m^jx\subset R
                       \text{ for some }j\},\qquad
 E=L\cap R^\sharp.
\]
The polynomial coaddition and comultiplication
tensors are uniquely determined if and only if
$E=R$, equivalently $\Omega(R,L)=R[X]$.
No assumption on normalization is needed.
\end{saresult}
\begin{proof}
The coefficient argument of B.034 gives
$\Omega(R,L)=\Omega(R,E)$, and $E=R$ is precisely
the polynomial-ring case, which is positive.
For necessity put $A=\widehat R$,
$Q=K\otimes_R A$ and $B=E\otimes_R A\subset Q$.
B.051 identifies the integer-valued modules and
their tensor maps. Every element of $B/A$ is
killed by a power of the maximal ideal of $A$.

There is a complete coefficient discrete valuation
subring $V\subset A$ to which B.052 applies.
In characteristic $p$ use the coefficient field
$k\subset A$ and $V=k[[t]]$, where
$0\ne t\in\mathfrak m\subset R$; flatness makes
$t$ a nonzerodivisor, so this power-series map
is injective. In mixed characteristic use the
Cohen coefficient ring with uniformizer $p$;
$p\in R$ is again a nonzerodivisor. In both
cases the uniformizer belongs to the
multiplicative set $W=R\setminus\{0\}$ defining
$Q=W^{-1}A$.

If the tensor map is injective, B.052 yields
$B=A+eQ$, where $eQ$ is finite-dimensional over
$F=\operatorname{Frac}(V)$ and is contained in
$A[1/t]$. If $e=0$, faithful flatness gives
$E=R$, as required. Suppose $e\ne0$ and choose
a local Artinian factor $Q_0=e_0Q$ of the
finite-dimensional algebra $eQ$.
The image $A_0=e_0A\subset Q_0$ is a quotient
of $A$, hence a complete Noetherian local ring,
and $A_0[1/t]=Q_0$.
It is finite over $V$: apply B.048 to this
submodule of the finite-dimensional $F$-space
$Q_0$. A vector-space summand would contradict
$t$-adic separatedness of the finite $A$-module
$A_0$. Thus $A_0$ is a one-dimensional complete
local order, with the same residue field as $A$.

We detail why the selected-factor arguments remain
valid although the whole ring $A$ need not be
finite over $V$. If $Q_0$ is nonreduced, apply
the coefficient-field and finite-order construction
in B.045 to the one-dimensional ring $A_0$.
It first chooses a new coefficient discrete valuation
ring $V_0\subset A_0$ and a coefficient field
$F_1\subset Q_0$ whose embedding agrees with
$\operatorname{Frac}(V_0)\subset Q_0$. This compatibility
is supplied by the separable primitive-element and
nilpotent Newton-lifting argument of B.045.
Let $V_1\subset F_1$ be the finite integral closure
of $V_0$, with uniformizer $u$. The ring
$O=V_1[A_0]\subset Q_0$ is finite over $V_1$,
contains $A_0$, and spans $Q_0$ over $F_1$.
It is also finite over $A_0$: a finite set of
$V_0$-module generators of $V_1$ generates $O$
as an $A_0$-module, because the two copies of
$V_0$ agree. Equivalently, these generators are
integral over $V_0\subset A_0$. Since $A_0$ is
finite over the original $V$, so is $O$.
No inclusion $V_1\subset A_0$ is asserted.
Since $Q_0\subset eQ\subset A[1/t]$, a single
power $t^N$ sends $e_0O$ into $A$.
Multiplication by $t$ is invertible on $Q_0$;
on the finite free $V_1$-module $O$ its cokernel
is torsion. Hence, for some $c$, one has
\[
                        u^ce_0O\subset t^Ne_0O\subset A.
                                                               \tag{B53.1}
\]
All the subsequent steps of B.045 use only the
projection $A\to A_0\subset O$, the inclusion
(B53.1), and nilpotent inputs supported on $Q_0$.
They give polynomials $F_r,H_r$ supported on
$Q_0$, hence with coefficients in
$e_0Q\subset B$, taking all of $A$ into $A$.
The very same pairing (B45.3), with scalar action
through $A\to A_0\to V_1$, is $A$-balanced and
detects their nonzero kernel tensor.
Thus a selected nonreduced factor is impossible.

If $Q_0$ is a field, its valuation ring is a
finite complete discrete valuation overring of
$A_0$. The same denominator-clearing argument
gives $u^ce_0V_1\subset A$ for some $c\ge1$.
One must have $e_0\ne1$. Otherwise $Q=Q_0$
would be finite-dimensional over $F$, and B.048
applied to the $t$-adically separated submodule
$A\subset Q=A[1/t]$ would make $A$ finite
over $V$, contradicting $\dim A=\dim R\ge2$.
The supported-input finite-difference proof of
B.050 now applies word for word: with
$x=u^ce_0$, every $h(x)-h(0)$ lies in the
proper ideal $A\cap e_0Q\subset\mathfrak m_A$,
while a supported congruence indicator $G$ has
$G(0)=0$ and valuation exactly $c$ at
$x^{c+1}$. Iterating multiplicative tensor
expressions would put this value in
$\mathfrak m_A^{c+1}$, a contradiction.
Thus a selected field factor is also impossible
if all multiplicative expressions exist.

Both possibilities contradict the joint
injectivity-and-existence hypothesis. Hence
$e=0$, $B=A$, and $E=R$ by faithful flatness.
This proves necessity and the theorem.
\end{proof}

\begin{saresult}{B.054}{Proved}{Complete Noetherian-domain classification with all torsion-free coefficients}
Let $R$ be a nonfield Noetherian domain with
fraction field $K$, and let $R\subset S$ be an
$R$-torsion-free unital extension. Put
$L=S_{\mathrm{rat}}$.
The polynomial coaddition and comultiplication
tensors are uniquely determined exactly under
the following conditions, imposed only at
maximal ideals $\mathfrak m$ with finite residue field:
\[
\begin{gathered}
 \dim R_{\mathfrak m}=1:\quad
 L_{\mathfrak m}=R_{\mathfrak m}\quad\hbox{or}\\
 \bigl(L_{\mathfrak m}=K\text{ and }
        \widehat{R_{\mathfrak m}}\text{ is a domain}\bigr);\\[2mm]
 \dim R_{\mathfrak m}\ge2:\quad
 L_{\mathfrak m}\cap R_{\mathfrak m}^{\sharp}=R_{\mathfrak m}.
\end{gathered}                                                   \tag{B54.1}
\]
There is no finite-normalization or finite bound on
Krull dimension. In every positive case all
multivariable tensor maps are isomorphisms and
the compatible plethory is unique.
For $S=K$, the first condition says finite local
normalization with one branch, and the second
is exactly $\operatorname{depth}R_{\mathfrak m}\ge2$.
\end{saresult}
\begin{proof}
Use the rational-coefficient reduction B.014 and
the localization theorem B.033. Infinite-residue
local rings give the ordinary polynomial ring.
At finite-residue maximal ideals, B.050 gives
the complete dimension-one criterion and B.053
gives the complete higher-dimensional criterion.
In the higher-dimensional positive case the
module is the polynomial ring, so every tensor
map is an isomorphism. In dimension one this
was proved in B.050. Test kernels and cokernels
at every maximal ideal and apply B.009 for
the compatible plethory. Finally B.046 identifies
the dimension-one completion condition with
finite local normalization, and B.034 identifies
$R_{\mathfrak m}^{\sharp}=R_{\mathfrak m}$ with
depth at least two. These arguments are local
and do not require a bound on the dimensions
of all maximal localizations simultaneously.
\end{proof}
